\documentclass[english, parskip=full]{scrartcl}
\usepackage[english]{babel}  % Sprache auf Deutsch 
\usepackage[utf8]{inputenc}  % Kodierung der Datei
\usepackage[left=2cm, right=2cm, top=2cm, bottom=3cm, footskip=1.5cm]{geometry}
\usepackage{color}
\usepackage{graphicx, gensymb}
\usepackage{amsfonts, cancel, amssymb}
\usepackage{amsmath, amsthm, array, esint}
\usepackage{bm} % bold math symbols, e.g. \bm{\sigma} etc.
\usepackage{booktabs, multirow}  % schönere Tabellen
%\usepackage{scrlayer-scrpage}    % Manipulation des Seitenstils
%\pagestyle{scrheadings}  % Stil für die Seite setzen
\usepackage{tabularx}
\usepackage{setspace}
%\automark{section}  % Abschnittsnamen als Seitenbeschriftung 
%\setheadsepline{.5pt}    % Kopzeile durch Linie abtrennen
\usepackage[hidelinks]{hyperref}  % Links und weitere PDF-Features
\usepackage{typearea, tikz, pgffor, verbatim}
\usetikzlibrary{calc, arrows.meta,arrows, graphs, quotes, positioning, angles, babel}
\usepackage[cache=false, outputdir=build]{minted}
\usemintedstyle{vs}
\usepackage{placeins} % provides \FloatBarrier

\definecolor{bg}{rgb}{0.9,0.9,0.9}
\newcommand\numberthis{\addtocounter{equation}{1}\tag{\theequation}}
\usepackage{svg}
\usepackage{siunitx}
\usepackage{physics}
\usepackage{tensor}
\usepackage{slashed}
\usepackage{bbm}
\usepackage[compat=1.1.0]{tikz-feynman}

\usepackage{caption}
\captionsetup{
    % width=.8\textwidth,
    % indention=1cm,
    margin=0.5cm,
    format=plain,
    labelfont=bf
}
%\usepackage{subcaption}
%\subcaptionsetup{
%    indention=0cm,
%    format=hang,
%    margin=0.5cm,
%    labelfont=normalfont
%}
\usepackage[english,capitalise]{cleveref}
\usepackage[
    backend=biber,
    sorting=none,
    style=phys,
    giveninits=true,
    sortcites=true,
    citestyle=numeric-comp,
    % url=false,
    % doi=false,
    biblabel=brackets,
    % eprint=false,
    maxcitenames=3]{biblatex}
\usepackage{csquotes}
%\addbibresource{bibliography.bib}

\usepackage{mathtools}
% use \abs for absolute values
% use \abs for \left ... \right version and \abs[\bigg for \bigg version etc.
\let\abs\undefined
\let\bra\undefined
\let\braket\undefined
\DeclarePairedDelimiter\abs{\vert}{\vert}
\DeclarePairedDelimiter\kla{(}{)}
\DeclarePairedDelimiter\klb{[}{]}
\DeclarePairedDelimiter\klc{\{}{\}}
\DeclarePairedDelimiter\bra{\langle}{\rangle}
\DeclarePairedDelimiterX\brb[2]{\langle #1}{\rangle}{\delimsize\vert #2 }
\DeclarePairedDelimiterX\brc[3]{\langle #1}{\rangle}{\delimsize\vert #2 \delimsize\vert #3 }

\renewcommand{\i}{\text{i}}
\newcommand{\e}{\text{e}}
\DeclareMathOperator{\sgn}{sgn}
\newcommand{\kom}[1]{\overset{\substack{#1 \\ |}}{=}}
\newcommand{\koml}[2]{\overset{\substack{#1 \\ |}}{#2}}
\newcommand{\intx}[4]{\int_{#1}^{#2} #3 \,\mathrm{d}#4}
\newcommand{\intr}[2]{\int_{\mathbb{R}} #1 \, \mathrm{d}#2}
\newcommand{\intrnd}[3]{\int_{\mathbb{R}^{#1}} #2 \, \mathrm{d}^{#1}#3}
\newcommand{\intrpi}[2]{\int_{\mathbb{R}} #1 \, \frac{\mathrm{d}#2}{2\pi}}
\newcommand{\intrndpi}[3]{\int_{\mathbb{R}^{#1}} #2 \, \frac{\mathrm{d}^{#1}#3}{(2\pi)^{#1}}}
\newcommand{\oper}[2]{\vert #1 \rangle \langle #2 \vert}
\newcommand{\Text}[1]{\text{#1}}    % this is relevant because \text does not autocomplete to the default '\text{@}' but rather '\textbf' etc.
\newcommand{\powe}[1]{\e^{#1}}
\newcommand{\pow}[2]{{#1}^{#2}}
\newcommand{\Ket}[1]{\vert #1 \rangle}
\newcommand{\Bra}[1]{\langle #1 \vert}
\newcommand*{\Fourier}{\textsc{Fourier}\ }
\newcommand*{\F}{\mathcal{F}}
\DeclareMathOperator{\arsinh}{arsinh}
\DeclareMathOperator{\arcosh}{arcosh}
\DeclareMathOperator{\artanh}{artanh}
\DeclareMathOperator{\sinc}{sinc}
\DeclareMathOperator{\cscc}{cscc}
\DeclareMathOperator{\sinhc}{sinhc}
\DeclareMathOperator{\cschc}{cschc}
\newcommand{\conv}{\mathop{\scalebox{3.5}{\raisebox{-0.38ex}{\makebox[0.5\width][c]{$\ast$}}}}\limits}
\newcommand*{\unity}{\text{\usefont{U}{bbold}{m}{n}1}}   % operator one from :: https://tex.stackexchange.com/questions/566780/import-mathbb1-character-from-the-unicode-math-package

\renewcommand{\d}{\text{d}}
\renewcommand{\r}{\text{r}}
\renewcommand{\t}{\text{t}}
\renewcommand{\a}{\text{a}}
\renewcommand{\o}{\text{o}}
\renewcommand{\F}{\text{F}}
\renewcommand{\c}{\text{c}}
\newcommand{\A}{\text{A}}
\newcommand{\B}{\text{B}}
\newcommand{\R}{\text{R}}
\newcommand{\M}{\text{M}}
\newcommand{\m}{\text{m}}
\newcommand{\s}{\text{s}}
\newcommand{\f}{\text{f}}
\newcommand{\K}{\text{K}}
\newcommand{\hc}{\text{h.c.}}
\newcommand{\kf}{k_\F}
\newcommand{\vf}{v_\F}
\newcommand{\const}{\text{const.}}
\renewcommand{\L}{\text{L}}
\newcommand{\gray}[1]{\bgroup\color{white!40!black}#1\egroup}
\newcommand{\TODO}[1]{\bgroup\color{red!60!black}#1\egroup}
\newcommand\QnA[1]{\bgroup\color{blue}#1\egroup}
\newcommand\highlight[1]{\bgroup\color{green!60!black}#1\egroup}

\newcommand{\cabx}[3]{\begin{smallmatrix}\gray{A}&\gray{:}&#1 \\ \gray{B}&\gray{:}&#2 \\ \gray{\Xi}&\gray{:}&#3\end{smallmatrix}}
\newcommand{\zabx}[3]{\begin{smallmatrix}\gray{A}&\gray{:}&#1 \\ \gray{B}&\gray{:}&#2 \\ \gray{Z}&\gray{:}&#3\end{smallmatrix}}
\newcommand{\abx}[3]{\begin{smallmatrix}#1 \\ #2 \\ #3\end{smallmatrix}}

\newcommand{\normord}[1]{
  {:\mathrel{\mspace{1mu}#1\mspace{1mu}}:}
}
\newcommand{\varnormord}[1]{
  {\mathrel{\mspace{1mu}\substack{\ast\\[-1.5pt] \ast}#1\substack{\ast\\[-1.5pt] \ast}\mspace{1mu}}}
}

% punctuation styles (see https://tex.stackexchange.com/questions/56063/how-to-properly-typeset-footnotes-superscripts-after-punctuation-marks)
\let\origfootnote\footnote
\renewcommand{\footnote}[1]{\kern.06em\origfootnote{#1}}
\newcommand{\punctfootnote}[1]{\kern-.06em\origfootnote{#1}}

% Multiline equations
\newcommand{\bsplit}[1]{\begin{aligned}[t]#1\end{aligned}}

\newcommand*{\titel}{One-dimensional chain with 2-atomic basis and linear dispersion}
\newcommand*{\autor}{Joachim Schwardt}

\hypersetup{pdfauthor={\autor}, pdftitle={\titel}}

%\titlehead{titlehead}
\subject{Quantum Physics in 1D}
\title{\titel}
\author{\autor}
\date{\today}

\begin{document}

%\begin{titlepage}
\maketitle  % Titel setzen
%\tableofcontents  % Inhaltsverzeichnis setzen
%\end{titlepage}

\section{Model Hamiltonian}
We consider the model $H=H_0 + H_U$ where
\begin{align}
H_0 &= \vf\sum_k k \varnormord{c_{k,\A}^\dagger c_{k,\B}} + \hc \equiv \vf\sum_k \varnormord{\Psi_k^\dagger k\sigma_x \Psi_k},\\
H_U &= \sum_{j,\kappa = \A,\B} U_\kappa n_{j,\kappa} n_{j+1,\kappa}.
\end{align}
Here $U_\kappa\in\mathbb{R}$, $\Psi = (c_{k,\A}, c_{k,\B})^T$ and we set $\hbar=1$.
Note that we have two different kinds of normal-ordering.
The one used above can readily be defined by
\begin{align}
\varnormord{A} &= A - \bra*{A}
\end{align}
for any operator $A$.
It has the interpretation of simply removing the vacuum expectation value.
The other normal-ordering requires a basis of annihilation and creation operators $b$ and $b^\dagger$ that make up the operator $A$.
Then $\normord{(\dots)}$ reorders products of $b$ and $b^\dagger$ such that all $b$ appear on the right of the expression, and all $b^\dagger$ to appear on the left.
While this is conceptually simple, it is not all that simple to write it down in a formal manner.
Instead, just consider the example $\normord{b^\dagger b} = b^\dagger b$, $\normord{bb^\dagger}=b^\dagger b$ and $\normord{\kla*{b + bb^\dagger b}} = b + b^\dagger b^2$.
Note that in the special case of a bilinear operator the two normal-ordering prescriptions may actually have the same meaning, since $\varnormord{b^\dagger b} = b^\dagger b - \brc*{0}{b^\dagger b}{0} = b^\dagger b$ due to $b\Ket{0} = 0$.
However, in general this is not correct and we will therefore distinguish between both ideas.\\
For now we want to treat $H_0$ as the ``exactly solvable'' part of the model and $H_U$ as the perturbation.
This requires diagonalizing $\sigma_x$ by
\begin{align}
\sigma_x &= U \sigma_z U^\dagger
\end{align}
with the transformation matrix $U = \frac{1}{\sqrt{2}} \begin{pmatrix}
1 & 1 \\ 1 & -1
\end{pmatrix}$.
We may use this to change into the $\ell=\L,\R$-basis 
\begin{align}
\Psi &\rightarrow U\Psi = \frac{1}{\sqrt{2}} \begin{pmatrix}
c_\A + c_\B \\ c_\A - c_\B
\end{pmatrix} \equiv \begin{pmatrix}
c_\R \\ c_\L
\end{pmatrix}.
\end{align}
In this basis we find
\begin{align}
H_0 &= \vf\sum_{k,\ell} k\ell n_{k\ell}
\end{align}
where $\ell = \L$ will be treated as $-1$ and $\ell = \R$ as $+1$.
Applying the inverse Fourier transform
\begin{align}
c_{k,\ell} &= \frac{1}{\sqrt{L}} \intr{\powe{-\i kx} c_{x,\ell}}{x}
\end{align}
allows us to rewrite $H_0$ as\footnote{Note that $\frac{1}{L}\sum_k \powe{\i kx} \simeq \delta(x)$ for $k=\frac{2\pi}{L}n$ with $n\in\mathbb{Z}$ and large $L$.}
\begin{align}
H_0 &= \vf\sum_{k,\ell} k\ell \varnormord{\intr{\intr{\powe{\i kx} c_{x,\ell}^\dagger \powe{-\i kx'} c_{x',\ell}}{x}}{x'}} \\
&= \vf\sum_\ell \ell \varnormord{\intr{c_{x,\ell}^\dagger \intr{c_{x',\ell} \i \partial_{x'} \frac{1}{L} \sum_k \powe{\i k (x - x')} }{x'}}{x}} \\
&= \vf\sum_\ell \ell \intr{\varnormord{c_{x,\ell}^\dagger (-\i\partial_x) c_{x,\ell}}}{x}.
\end{align}
We will use one of the standard bosonization prescriptions
\begin{align}
c_{x,\ell} &\rightarrow \sqrt{a} \psi_\ell(x) = \sqrt{a} \frac{\eta_\ell}{\sqrt{2\pi \alpha}} \powe{-\i\varphi_\ell(x)} \label{eqn:bosonization_prescription}
\end{align}
with the short-distance cut-off $\alpha$ and a lattice spacing $a$.
For distinct species $\ell\neq\ell'$ the Klein factors satisfy the usual fermionic anti-commutation relations $\klc*{\eta_\ell, \eta_{\ell'}} = 0$ and $\klc[\big]{\eta_\ell, \eta_{\ell'}^\dagger} = 0$.
In the thermodynamic limit the Klein factors obey the Majorana algebra, and therefore satisfy $\eta_\ell^\dagger \eta_\ell \equiv 1 \equiv \eta_\ell\eta_\ell^\dagger$.
Furthermore $\eta_\ell^\dagger \equiv \eta_\ell$, which allows a representation in terms of Pauli matrices.
To summarize, for $L\rightarrow\infty$ the Klein factors are hermitian, unitary operators that square to unity and anti-commute for distinct species.\\
Note that the chiral fields $\varphi_\ell = \ell\phi - \theta$ can be split into a part containing creation operators and a part containing annihilation operators via $\varphi_\ell = \varphi_\ell^{(+)} + \varphi_\ell^{(-)}$.
These parts satisfy the commutation relation
\begin{align}
\klb*{\varphi_\ell^{(+)}(x), \varphi_{\ell'}^{(-)}(x')} &= \delta_{\ell\ell'} \log\klb*{\frac{2\pi}{L} \kla*{\alpha + \i\ell (x-x')} } \label{eqn:commutator_varphi_ell_plus_x_varphi_ellp_minus_xp},
\end{align}
which can be derived by expressing $\varphi_\ell$ in terms of bosonic creation and annihilation operators $b_q$ and $b_q^\dagger$.\punctfootnote{See \texttt{bosonization{\_}dictionary.tex} and references therein.}
We will also make use of
\begin{align}
\klb*{\varphi_\ell(x), \varphi_{\ell'}(x')} &= \sum_{pp'} \klb*{\varphi_\ell^{(p)}(x), \varphi_{\ell'}^{(p')}(x')} = \sum_p \delta_{\ell\ell'} p \log\klb*{\frac{2\pi}{L} \kla*{\alpha + p\i\ell (x-x')} } \\
&= \delta_{\ell\ell'} \log \klb*{ \frac{\alpha + \i \ell (x-x')}{\alpha - \i \ell (x-x')}} \label{eqn:commutator_varphi_ell_x_varphi_ellp_xp},\\
\klb*{\phi^{(+)}(x), \phi^{(-)}(x')} &= \frac{1}{4} \sum_{\ell\ell'} \klb*{\ell\varphi_\ell^{(+)}(x), \ell'  \varphi_{\ell'}^{(-)}(x')} = \frac{1}{4} \sum_\ell \log \klb*{\frac{2\pi}{L} \kla*{\alpha + \i\ell (x-x')} } \\
&= \frac{1}{2} \log \klb*{\frac{2\pi}{L} \sqrt{\alpha^2 + (x-x')^2}} \label{eqn:commutator_phi_plus_x_phi_minus_xp}.
\end{align}
Since the commutator is $\mathbb{C}$-valued, Baker-Campbell-Hausdorff here reads $\powe{A}\powe{B} = \powe{A+B} \powe{\frac{1}{2} \klb*{A,B}}$.
Similarly, $\powe{A}\powe{B} = \powe{B}\powe{A} \powe{\klb*{A,B}}$.
We may then use these identities to evaluate
\begin{align}
\powe{\i m\varphi_\ell(x)} &= \powe{\i m\varphi_\ell^{(+)}(x)} \powe{\i m\varphi_\ell^{(-)}(x)} \powe{-\frac{1}{2} \klb*{\i m\varphi_\ell^{(+)}(x),\, \i m\varphi_\ell^{(-)}(x)}} = \normord{\powe{\i m\varphi_\ell(x)}} \kla*{\frac{2\pi \alpha}{L}}^{\frac{m^2}{2}}.
\end{align}
This derivation encompasses the essence of how to correctly treat products of these kinds of exponential or vertex operators.
In a similar way we may now also compute
\begin{align}
\psi_\ell^\dagger(x) \psi_\ell(x+\epsilon) &= \frac{\eta_\ell^\dagger \eta_\ell}{2\pi \alpha} \powe{\i\varphi_\ell(x)} \powe{-\i\varphi_\ell(x+\epsilon)} = \frac{1}{2\pi \alpha} \normord{\powe{\i\varphi_\ell(x)}} \normord{\powe{-\i\varphi_\ell(x+\epsilon)}} \frac{2\pi \alpha}{L} \\
&= \frac{1}{L} \powe{\i\varphi_\ell^{(+)}(x)} \powe{\i\varphi_\ell^{(-)}(x)} \powe{-\i\varphi_\ell^{(+)}(x+\epsilon)} \powe{-\i\varphi_\ell^{(-)}(x+\epsilon)} \\
&= \frac{1}{L} \powe{\i\varphi_\ell^{(+)}(x)} \powe{-\i\varphi_\ell^{(+)}(x+\epsilon)} \powe{\i\varphi_\ell^{(-)}(x)} \powe{-\i\varphi_\ell^{(-)}(x+\epsilon)} \klb*{\frac{2\pi}{L} \kla*{\alpha + \i\ell\epsilon}}^{-1} \\
&= \frac{1}{2\pi} \frac{1}{\alpha + \i\ell\epsilon} \normord{\powe{\i \varphi_\ell(x) - \i\varphi_\ell(x+\epsilon)}}.
\end{align}
In this expression we can safely take the limit $\alpha\rightarrow 0$, and due to the normal-ordering also Taylor-expand $\varphi_\ell(x+\epsilon)$ in the exponential.
Dropping irrelevant constants we then find
\begin{align}
\lim_{\epsilon\rightarrow 0^+} \varnormord{\psi_\ell^\dagger(x) \psi_\ell(x+\epsilon)} &= \lim_{\epsilon\rightarrow 0^+} \frac{-\i\ell}{2\pi\epsilon} \klb*{\normord{\powe{-\i \epsilon\partial_x\varphi_\ell(x)}} - \bra*{\normord{\powe{-\i \epsilon\partial_x\varphi_\ell(x)}}}} \\
&= \lim_{\epsilon\rightarrow 0^+} \frac{-\i\ell}{2\pi\epsilon} \normord{\klb[\big]{-\i\epsilon\partial_x\varphi_\ell(x)}} = \frac{-\ell}{2\pi} \partial_x\varphi_\ell(x) \equiv \rho_\ell
\end{align}
where normal-ordering can be dropped since $\varphi_\ell$ is given by a linear combination of $b_q$ and $b_q^\dagger$ and is thus already normal-ordered by construction.
To bosonize $H_0$ we consider the same calculation, but due to the derivative need to expand to a higher order in the point-splitting parameter $\epsilon$.
We have
\begin{align}
\varnormord{\psi_\ell^\dagger(x) \i\partial_x \psi_\ell(x+\epsilon)} &= \i\partial_\epsilon \varnormord{\psi_\ell^\dagger (x) \psi_\ell(x+\epsilon)} = \i\partial_\epsilon \frac{-\i\ell}{2\pi\epsilon} \klb*{\normord{\powe{-\i\epsilon\partial_x\varphi_\ell(x) - \i \frac{\epsilon^2}{2} \partial_x^2\varphi_\ell(x)}} - 1} \\
&= \i\partial_\epsilon \frac{-\i\ell}{2\pi\epsilon} \normord{\klb*{-\i\epsilon\partial_x\varphi_\ell(x) - \i \frac{\epsilon^2}{2} \partial_x^2\varphi_\ell(x) - \frac{\epsilon^2}{2} \kla*{\partial_x\varphi_\ell(x)}^2}}\\
&= \frac{-\ell}{4\pi} \normord{\klb*{\kla*{\partial_x\varphi_\ell(x)}^2 + \i \partial_x^2\varphi_\ell(x)}}.
\end{align}
We want to work in the continuum limit and therefore take the limit $a\rightarrow 0$ whenever possible.
The free Hamiltonian then becomes\footnote{As usual $\sum_j \rightarrow \int \,\frac{\d x}{a}$.}
\begin{align}
H_0 &= \vf\sum_\ell \ell \intr{\varnormord{\psi_\ell^\dagger(x) (-\i\partial_x) \psi_\ell(x)}}{x} = \frac{\vf}{2\pi} \intr{\varnormord{\klb*{\kla*{\partial_x\theta}^2 + \kla*{\partial_x\phi}^2}} }{x}
\end{align}
where the total derivative $\partial_x^2\varphi_\ell$ dropped out and we used $\sum_\ell \kla*{\partial_x\varphi_\ell}^2 = 2 \kla*{\partial_x\theta}^2 + 2\kla*{\partial_x\phi}^2$.
\section{Bosonization of the interaction}
Before we turn to the bosonization of the interaction Hamiltonian we first perform the basis transformation used to diagonalize the free Hamiltonian
\begin{align}
c_\kappa &= \frac{c_\R + \kappa c_\L}{\sqrt{2}}.
\end{align}
to get $n_\kappa = \frac{1}{2} \kla[\big]{c_\R^\dagger + \kappa c_\L^\dagger} \kla*{c_\R + \kappa c_\L} = \frac{1}{2} \kla*{n_\R + n_\L + \kappa c_\R^\dagger c_\L + \hc} = \frac{1}{2} \sum_\ell \kla[\big]{n_\ell + \kappa c_\ell^\dagger c_{-\ell}}$. 
Using this we find
\begin{align}
H_U &= \frac{1}{4} \sum_{j,\kappa=\A,\B} U_\kappa \sum_{\ell,\ell' = \R,\L} \kla*{n_{j,\ell} + \kappa c_{j,\ell}^\dagger c_{j,-\ell}} \kla*{n_{j+1,\ell'} + \kappa c_{j+1,\ell'}^\dagger c_{j+1,-\ell'}} = H_{U_+} +  H_{U_-}
\end{align}
where $U_\pm = \frac{U_\A \pm U_\B}{2}$ and
\begin{align}
H_{U_+} &= \frac{U_+}{2} \sum_{j,\ell,\ell'} \kla*{n_{j,\ell} n_{j+1,\ell'} + c_{j,\ell}^\dagger c_{j,-\ell} c_{j+1,\ell'}^\dagger c_{j+1,-\ell'}},\\
H_{U_-} &= \frac{U_-}{2} \sum_{j,\ell,\ell'} \kla*{c_{j,\ell}^\dagger c_{j,-\ell} n_{j+1,\ell'} + n_{j,\ell} c_{j+1,\ell'}^\dagger c_{j+1,-\ell'} }.
\end{align}
We will implicitly assume integration over $x$ and make use of the approximation $F(x+a) \simeq F(x)$ for the slowly-varying fields $F\in \klc*{\varphi_\ell, \phi, \theta}$ whenever applicable.\footnote{Note that this needs to be treated with caution whenever the exponent would completely vanish due to such an approximation, as in the previous $\psi_\ell^\dagger \psi_\ell$ example.} 
First, the density-density term simply renormalizes the free Luttinger parameters since
\begin{align}
\sum_{\ell,\ell'} \frac{-\ell}{2\pi} \partial_x\varphi_\ell \frac{-\ell'}{2\pi} \partial_x\varphi_{\ell'} = \frac{1}{4\pi^2} \sum_{\ell,\ell'} \ell \ell' \kla*{\ell \partial_x\phi - \partial_x\theta} \kla*{\ell' \partial_x\phi - \partial_x\theta} = \frac{1}{\pi^2} \kla*{\partial_x\phi}^2.
\end{align}
The second term behaves very differently for $\ell=\ell'$ and $\ell\neq\ell'$, so we consider those cases separately.
Note that the first vanishes in the continuum limit $\rightarrow 0$ because
\begin{align}
c_{j,\ell}^\dagger c_{j,-\ell} c_{j+1,\ell}^\dagger c_{j+1,-\ell} &= - c_{j,-\ell} c_{j,\ell}^\dagger c_{j+1,\ell}^\dagger c_{j+1,-\ell} \overset{a\,\rightarrow\,0}{\longrightarrow} c_{j,-\ell} \kla[\big]{c_{j,\ell}^\dagger}^2 c_{j,-\ell} \equiv 0.
\end{align}
This is consistent with the bosonization procedure, as shown in \autoref{subsec:correct_order_of_limits}.
The second case gives a more interesting result.
First note that
\begin{align}
c_{j,\ell}^\dagger c_{j,-\ell} c_{j+1,-\ell}^\dagger c_{j+1,\ell} &= c_{j,\ell}^\dagger c_{j+1,\ell} c_{j,-\ell} c_{j+1,-\ell}^\dagger = - c_{j,\ell}^\dagger c_{j+1,\ell} c_{j+1,-\ell}^\dagger c_{j,-\ell}.
\end{align}
Just as in the bosonization of the free Hamiltonian we now consider the terms
\begin{align}
\psi_\ell^\dagger(x) \psi_\ell(x+a) &= \frac{1}{L} \normord{\powe{\i\varphi_\ell(x)}} \normord{\powe{-\i\varphi_\ell(x+a)}} = \frac{1}{L} \normord{\powe{ \i\varphi_\ell(x) - \i\varphi_\ell(x+a)}} \powe{\klb*{\i\varphi_\ell^{(-)}(x), -\i\varphi_\ell^{(+)}(x+a)}} = \\
&= \frac{1}{2\pi} \frac{1}{\alpha +\i\ell a} \normord{\powe{-\i a\partial_x\varphi_\ell - \i \frac{a^2}{2} \partial_x^2\varphi_\ell + \mathcal{O}(a^3)}} = \\
&= \frac{-\i\ell }{2\pi a} \normord{\klb*{1 - \i a\partial_x\varphi_\ell - \i\frac{a^2}{2} \partial_x^2\varphi_\ell - \frac{a^2}{2}\kla*{\partial_x\varphi_\ell}^2 + \mathcal{O}(a^3)}} = \\
&= \frac{-\i\ell}{2\pi a} - \frac{\ell}{2\pi} \partial_x\varphi_\ell + \frac{\i\ell a}{4\pi} \klb*{\i\partial_x^2\varphi_\ell + \normord{\kla*{\partial_x\varphi_\ell}^2}} + \mathcal{O}(a^2),\\
\psi_{-\ell}^\dagger(x+a) \psi_{-\ell}(x) &= \frac{-\i\ell }{2\pi a} + \frac{\ell}{2\pi} \partial_x\varphi_{-\ell} + \frac{\i\ell a}{4\pi} \klb*{-\i\partial_x^2\varphi_{-\ell} + \normord{\kla*{\partial_x\varphi_{-\ell}}^2}} + \mathcal{O}(a^2).
\end{align}
To order $\mathcal{O}(a)$ the product $\psi_\ell^\dagger(x) \psi_\ell(x+a) \psi_{-\ell}^\dagger(x+a) \psi_{-\ell}(x)$ gives
\begin{align}
\frac{-1}{(2\pi a)^2} + \frac{-1}{4\pi^2} \partial_x\varphi_\ell \partial_x\varphi_{-\ell} + \frac{1}{8\pi^2} \klb*{-\i\partial_x^2\varphi_{-\ell} + \normord{\kla*{\partial_x\varphi_{-\ell}}^2} + \i\partial_x^2\varphi_\ell + \normord{\kla*{\partial_x\varphi_\ell}^2}}.
\end{align}
Dropping the constant and the total derivatives we get
\begin{align}
\frac{1}{4\pi^2} &\normord{\kla*{\ell \partial_x\phi - \partial_x\theta} \kla*{-\ell \partial_x\phi - \partial_x\theta}} + \frac{1}{4\pi^2} \normord{\klb*{\kla*{\partial_x\phi}^2 + \kla*{\partial_x\theta}^2 }} = \frac{1}{2\pi^2} \normord{\kla*{\partial_x\theta}^2}.
\end{align}
Summation on $\ell$ and the additional minus sign then give $\frac{-U_+a}{2\pi^2} \kla*{\partial_x\theta}^2$.
\TODO{TODO: Giamarchi (8.3) on p. 254 essentially corresponds to our interaction, and setting $\kf=0$ would completely eliminate any influence from that term.
The difference is that, in this section, he does not pay any attention to Klein factors whatsoever and takes a very different approach overall.
It seems unlikely that my derivation is more correct, but that the interaction term simply drops out when setting $\kf=0$ does not appear reasonable to me, so I am really not convinced by his result...}\\
We now turn to the third and fourth term of the interaction Hamiltonian.
Once again, before bosonizing we first manipulate the fermionic operators.
Starting with the case $\ell=\ell'$ we find
\begin{align}
c_{j,\ell}^\dagger c_{j,-\ell} n_{j+1,\ell} &= -c_{j,-\ell} c_{j,\ell}^\dagger c_{j+1,\ell}^\dagger c_{j+1,\ell} \overset{a\,\rightarrow \, 0}{\longrightarrow} -c_{j,-\ell} \kla[\big]{c_{j,\ell}^\dagger}^2 c_{j,\ell} \equiv 0,\\
n_{j,\ell} c_{j+1,\ell}^\dagger c_{j+1,-\ell} &= \kla*{1 - c_{j,\ell} c_{j,\ell}^\dagger} c_{j+1,\ell}^\dagger c_{j+1,-\ell} \overset{a\,\rightarrow \, 0}{\longrightarrow} c_{j,\ell}^\dagger c_{j,-\ell}.
\end{align}
Similarly for $\ell = -\ell'$ we find
\begin{align}
c_{j,\ell}^\dagger c_{j,-\ell} n_{j+1,-\ell} &= c_{j,\ell}^\dagger c_{j,-\ell} \kla*{1 - c_{j+1,-\ell} c_{j+1,-\ell}^\dagger} \overset{a\,\rightarrow \, 0}{\longrightarrow} c_{j,\ell}^\dagger c_{j,-\ell},\\
n_{j,\ell} c_{j+1,-\ell}^\dagger c_{j+1,\ell} &= -c_{j,\ell}^\dagger c_{j,\ell} c_{j+1,\ell} c_{j+1,-\ell}^\dagger \overset{a\,\rightarrow \, 0}{\longrightarrow} 0.
\end{align}
Thus we find that the interaction actually reduces to the quadratic expression 
\begin{align}
H_{U_-} &= U_- \sum_{j,\ell} c_{j,\ell}^\dagger c_{j,-\ell}
\end{align}
in the continuum limit.
Bosonizing this is rather straight-forward, since the chiral fields commute for different $\ell$.
Thus we find
\begin{align}
\frac{\eta_\ell^\dagger \eta_{-\ell}}{2\pi\alpha} \powe{\i\varphi_\ell - \i\varphi_{-\ell}} &= \frac{\eta_\ell^\dagger \eta_{-\ell}}{2\pi\alpha} \powe{2\i\ell \phi}
\end{align}
and in the thermodynamic limit with $\eta^\dagger=\eta$ summing on $\ell$ yields
\begin{align}
\frac{\eta_\R \eta_\L}{2\pi\alpha} 2\i\sin(2\phi).
\end{align}
Collecting terms our Hamiltonian now takes the form $H=H_0+H_I$ with
\begin{align}
H_0 &= \frac{\vf}{2\pi} \intr{\normord{\klb*{\kla*{1 - \frac{U_+a}{\pi\vf}} \kla*{\partial_x\theta}^2 + \kla*{1 + \frac{U_+a}{\pi\vf}} \kla*{\partial_x\phi}^2  }}}{x},\\
H_I &= \frac{U_- \eta_\R\eta_\L}{2\pi\alpha} \intr{2\i\sin(2\phi)}{x}
\end{align}
or by introducing the effective parameters $u = \vf\sqrt{1 - \frac{U_+^2a^2}{\pi^2\vf^2}}$ and $K = \sqrt{\frac{\pi\vf - U_+a}{\pi\vf + U_+a}}$
\begin{align}
H_0 &= \frac{u}{2\pi} \intr{\normord{\klb*{K\kla*{\partial_x\theta}^2 + \frac{1}{K} \kla*{\partial_x\phi}^2}}}{x},\\
H_I &= \frac{U_-}{2\pi\alpha} \sum_\ell \eta_{\ell}\eta_{-\ell} \intr{\powe{2\i\ell \phi}}{x}.
\end{align}
\subsection{On the correct order of limits}\label{subsec:correct_order_of_limits}
The bosonization procedure is, somewhat surprisingly, not very straightforward.
To understand why the naive approach of plugging the bosonization identity into the interaction Hamiltonian does not yield the correct result, one has to consider the order of limits.
Since we ultimately want to take $\alpha,a,\epsilon \rightarrow 0$ there are a priori six possible options with regard to the order of those limits.
For illustrative purposes we recalculate some of the previous results obtained by reordering the fermionic operators, but this time using only the bosonized fields.
First, note that since the chiral fields commute for different $\ell$ we only need to consider strings of three $c$'s, which share the same index.
For example all the complexity of $c_{j,\ell}^\dagger c_{j,-\ell} n_{j+1,\ell}$ is contained in $c_{j,\ell}^\dagger n_{j+1,\ell}$.
The latter is
\begin{align*}
\frac{\eta_\ell^\dagger \eta_\ell^\dagger \eta_\ell }{L^{3/2}} & \normord{\powe{\i\varphi_\ell(x)}} \normord{\powe{\i\varphi_\ell(x+a)}} \normord{\powe{-\i\varphi_\ell(x+a+\epsilon)}}  \\
&= \frac{\eta_\ell^\dagger }{L^{3/2}} \normord{\powe{\i\varphi_\ell(x)} \powe{\i\varphi_\ell(x+a)}} \frac{2\pi}{L} (\alpha + \i\ell a)  \normord{\powe{-\i\varphi_\ell(x+a+\epsilon)}}  \\
&= \frac{\eta_\ell^\dagger }{L^{3/2}} \normord{\powe{\i\varphi_\ell(x)} \powe{\i\varphi_\ell(x+a)} \powe{-\i\varphi_\ell(x+a+\epsilon)}} \frac{2\pi}{L} (\alpha + \i\ell a) \frac{L}{2\pi} \frac{1}{\alpha + \i\ell\epsilon}  \frac{L}{2\pi} \frac{1}{\alpha + \i\ell(\epsilon + a)} .\numberthis \label{eqn:c_ell_dagger_n_ell_bosonized_limit_order}
\end{align*}
Similarly $c_{j,\ell} n_{j+1,\ell}$ becomes
\begin{align*}
\frac{\eta_\ell \eta_\ell^\dagger \eta_\ell }{L^{3/2}} & \normord{\powe{-\i\varphi_\ell(x)}} \normord{\powe{\i\varphi_\ell(x+a)}} \normord{\powe{-\i\varphi_\ell(x+a+\epsilon)}}  \\
&= \frac{\eta_\ell }{L^{3/2}} \normord{\powe{-\i\varphi_\ell(x)} \powe{\i\varphi_\ell(x+a)}} \frac{L}{2\pi} \frac{1}{\alpha + \i\ell a} \normord{\powe{-\i\varphi_\ell(x+a+\epsilon)}}  \\
&= \frac{\eta_\ell }{L^{3/2}} \normord{\powe{-\i\varphi_\ell(x)} \powe{\i\varphi_\ell(x+a)} \powe{-\i\varphi_\ell(x+a+\epsilon)}}  \frac{L}{2\pi} \frac{1}{\alpha + \i\ell a} \frac{L}{2\pi} \frac{1}{\alpha + \i\ell\epsilon}  \frac{2\pi}{L} \kla[\big]{\alpha + \i\ell(\epsilon + a)}.\numberthis \label{eqn:c_ell_n_ell_bosonized_limit_order}
\end{align*}
As a last example consider $c_{j,\ell} c_{j+1,\ell} c_{j+1,\ell}^\dagger$ which becomes
\begin{align*}
\frac{\eta_\ell \eta_\ell \eta_\ell^\dagger }{L^{3/2}} & \normord{\powe{-\i\varphi_\ell(x)}} \normord{\powe{-\i\varphi_\ell(x+a)}} \normord{\powe{\i\varphi_\ell(x+a+\epsilon)}}  \\
&= \frac{\eta_\ell }{L^{3/2}} \normord{\powe{-\i\varphi_\ell(x)} \powe{-\i\varphi_\ell(x+a)}} \frac{2\pi}{L} (\alpha + \i\ell a) \normord{\powe{\i\varphi_\ell(x+a+\epsilon)}}  \\
&= \frac{\eta_\ell }{L^{3/2}} \normord{\powe{-\i\varphi_\ell(x)} \powe{-\i\varphi_\ell(x+a)} \powe{\i\varphi_\ell(x+a+\epsilon)}} \frac{2\pi}{L} (\alpha + \i\ell a) \frac{L}{2\pi} \frac{1}{\alpha + \i\ell\epsilon}  \frac{L}{2\pi} \frac{1}{\alpha + \i\ell(\epsilon + a)}. \numberthis \label{eqn:c_ell_c_ell_c_ell_dagger_bosonized_limit_order}
\end{align*}
For both of these examples the six different orders of limits are listed in \autoref{table:order_of_limits_for_c_ell_dagger_n_ell}.
Note that one should take $\epsilon\rightarrow 0$ as the last limit and that the second example always gives a singular result.
{\setlength{\tabcolsep}{15pt}
\begin{table}[!ht]
\centering
\begin{tabular}{l*{6}{c}}\toprule
 & \multicolumn{2}{c}{$a\rightarrow 0$} & \multicolumn{2}{c}{$\alpha\rightarrow 0$} & \multicolumn{2}{c}{$\epsilon\rightarrow 0$} \\ \midrule
 & $\alpha\rightarrow 0$ & $\epsilon\rightarrow 0$ & $a\rightarrow 0$ & $\epsilon\rightarrow 0$ & $a\rightarrow 0$ & $\alpha\rightarrow 0$ \\ \midrule
\autoref{eqn:c_ell_dagger_n_ell_bosonized_limit_order}: & 0 & $\frac{1}{\alpha}$ & 0 & $\frac{1}{\epsilon}$ & $\frac{1}{\alpha}$ & $\frac{1}{\alpha}$ \\ \midrule
\autoref{eqn:c_ell_n_ell_bosonized_limit_order}: & $\frac{1}{\alpha}$ & $\frac{1}{\alpha}$ & $\frac{1}{a}$ & $\frac{1}{\epsilon}$ & $\frac{1}{\alpha}$ & $\frac{1}{\alpha}$ \\ \midrule
\autoref{eqn:c_ell_c_ell_c_ell_dagger_bosonized_limit_order}: & 0 & $\frac{1}{\alpha}$ & 0 & $\frac{1}{\epsilon}$ & $\frac{1}{\alpha}$ & $\frac{1}{\alpha}$ \\ \bottomrule
\end{tabular}
\caption{Singular parts of for the six possible orders of limits for three examples of the bosonized expressions.
Note that most of the options for the first and third example lead to ambiguous results, and that one steers clear of difficulties so long as $\epsilon\rightarrow 0$ is the last limit to be taken.
Furthermore, the second example only produces ambiguous results and such terms should therefore be avoided by means of the fermionic commutation relations before bosonization.}
\label{table:order_of_limits_for_c_ell_dagger_n_ell}
\end{table}
}
In essence, difficulties arise because the anti-commutator $\klc*{c_{j,\ell}, c_{j,\ell}^\dagger} = 1$ is not so easy to replicate with the bosonization identity.
Careful analysis shows that $\klc*{\psi_\ell(x), \psi_\ell^\dagger(x+\epsilon)} = \delta(\epsilon)$, but since we are not integrating over $\epsilon$ this is a rather dangerous expression and most likely responsible for the difficulties.
\section{Single-particle Green's function in perturbation theory}
Our model turned out to be almost identical to the classical sine-Gordon model.
It is said to be integrable, but thus far the problem of obtaining the exact two-point correlation functions has resisted many attempts.
Here we want to take a more modest approach and merely treat the term in the first few orders of a perturbation series.
Recall that given some Hamiltonian $H=H_0+H_\Text{int}$ we calculate expectation values via
\begin{align}
\bra*{\mathcal{O}} &= \frac{1}{\mathcal{Z}} \int \mathcal{O} \powe{\i S[\Phi]}\,\mathcal{D}\Phi
\end{align} 
where $\mathcal{Z} = \int \powe{\i S[\Phi]}\,\mathcal{D}\Phi$ is the full partition function, $\mathcal{O}$ is some operator and $\Phi$ represents the fields appearing in the action.
Since the action decomposes into  a free and an interacting part via $S=S_0 + S_\Text{int}$ we find
\begin{align}
\mathcal{Z} &= \int \powe{\i S_0[\Phi]} \powe{\i S_\Text{int}[\Phi]}\,\mathcal{D}\Phi = \mathcal{Z}_0 \bra*{\powe{\i S_\Text{int}[\Phi]}}_0
\end{align}
where $\bra*{\mathcal{O}}_0$ denotes the average with respect to $S_0$.
Similarly
\begin{align}
\bra*{\mathcal{O}} &= \frac{1}{\mathcal{Z}}\int \mathcal{O} \powe{\i S_\Text{int}[\Phi]} \powe{\i S_0[\Phi]} \,\mathcal{D}\Phi = \frac{\mathcal{Z}_0 \bra*{\mathcal{O}\powe{\i S_\Text{int}}}_0}{\mathcal{Z}_0 \bra*{\powe{\i S_\Text{int}}}_0} = \frac{1}{\bra*{\powe{\i S_\Text{int}}}_0} \bra*{\mathcal{O}\powe{\i S_\Text{int}}}_0.
\end{align}
Expanding into a power series then yields
\begin{align}
\bra*{\mathcal{O}} &= \frac{1}{1 + \i \bra*{S_\Text{int}}_0 + \frac{\i^2}{2} \bra*{S_\Text{int}^2}_0 + \dots} \klc*{\bra*{\mathcal{O}}_0 + \i\bra*{\mathcal{O}S_\Text{int}}_0 + \frac{\i^2}{2} \bra*{\mathcal{O} S_\Text{int}^2} + \dots}\\
&= \bra*{\mathcal{O}}_0 + \i\bra*{\mathcal{O}S_\Text{int}}_0^c + \frac{\i^2}{2} \bra*{\mathcal{O}S_\Text{int}^2}_0^c - \i^2 \bra*{\mathcal{O}S_\Text{int}}_0^c \bra*{S_\Text{int}}_0 + \dots
\end{align}
where $\bra*{AB}^c = \bra*{AB} - \bra*{A}\bra*{B}$ denotes the connected average.
While we have converted the calculation of a single expectation value into a somewhat complicated series expansion, we only ever evaluated averages with respect to the free action $S_0$.
In our context of bosonization and the single-particle Green's function we can even go one step further and write down an explicit formula for the most general kind of expectation value we will encounter.
Assuming that $\phi(x,t)$ and $\theta(x,t)$ have $\mathbb{C}$-valued commutators for any $x$ and $t$ and that both $\bra*{\phi(0,0)^2}_0$ and $\bra*{\theta(0,0)^2}_0$ diverge one has
\begin{align}
\bra[\Big]{\prod_j \powe{\i A_j\phi(x_j,t_j) + \i B_j\theta(x_j,t_j)}} &= \powe{\frac{1}{2} \sum_{j<k} \klb*{\kla*{K A_jA_k + \frac{1}{K} B_jB_k} F_1(x_j-x_k,t_j-t_k) - \kla*{A_jB_k + A_kB_j} F_2(x_j-x_k,t_j-t_k)}}
\end{align}
under the neutrality condition $\sum_jA_j=0=\sum_jB_j$.
If this condition is not met, the expectation value gives zero.
Note that we can express all quantities and integrals in terms of $\xi_\pm$ via
\begin{align}
S_\Text{int} &= \frac{U_-}{2\pi\alpha u} \sum_\ell \eta_\ell \eta_{-\ell} \intrnd{2}{\powe{2\i\ell\phi}}{\xi},\\
\psi_\ell &= \frac{\eta_\ell}{\sqrt{2\pi\alpha}} \powe{-\i\varphi_\ell} = \frac{\eta_\ell}{\sqrt{2\pi\alpha}} \powe{-\i \ell \phi + \i \theta},\\
F(\xi) &= \log\klb*{\frac{\i\beta u}{\pi\alpha}} - \log\klb*{\csch\kla*{\frac{\pi \xi}{\beta u}}},\\
F_1(\xi) &= \log\klb*{
\frac{\i\beta u}{\pi\alpha}} - \frac{1}{2}\log\klb*{ \csch\kla*{\frac{\pi \xi_-}{\beta u}} \csch\kla*{\frac{\pi \xi_+}{\beta u}}}.
\end{align}
The general formula for the expectation value admits a recursion relation.
Writing $V_j = \powe{\i A_j\phi(\xi_j) + \i B_j \theta(\xi_j)}$ and $\Sigma_{jk} = \kla*{KA_jA_k + \frac{1}{K} B_jB_k}F_1(\xi_j-\xi_k) - \kla*{A_jB_k + A_kB_j}F_2(\xi_j-\xi_k)$ we have
\begin{align}
\bra[\Big]{\prod_{j=1}^{N+1} V_j } &= \powe{\frac{1}{2} \sum_{k=2}^{N+1} \sum_{j=1}^{k-1} \Sigma_{jk}} = \powe{\frac{1}{2} \sum_{k=2}^N \sum_{j=1}^{k-1} \Sigma_{jk} + \frac{1}{2}\sum_{j=1}^N \Sigma_{j,N+1} } = \bra[\Big]{\prod_{j=1}^NV_j} \powe{\frac{1}{2} \sum_{j=1}^N \Sigma_{j,N+1}}.
\end{align}
Another interesting observation is that swapping $V_1$ and $V_2$ only leads to the sign change $\xi_1 - \xi_2 \rightarrow \xi_2 - \xi_1$, leaving everything else unchanged.
Since $F^\ast(-\xi)=F(\xi)$ and $F_1^\ast(-\xi) = F_1(\xi)$ it follows that
\begin{align}
\bra[\Big]{\klc*{V_1, V_2} \prod_{j=3}^NV_j} &= 2\Re \klc*{\powe{\frac{1}{2} \Sigma_{12}}} \powe{\frac{1}{2} \sum_{j<k} \kla*{1 - \delta_{j1} \delta_{k2}} \Sigma_{jk}}.
\end{align}
Using these tools we will calculate the single-particle Green's function, i.e. we set $\mathcal{O} = \klc*{\psi_\ell(\xi), \psi_{\ell'}(0)^\dagger}$.
Since the perturbation only contains $\phi$ we can simply our general formula a bit\footnote{$A_{1} = -\ell$, $B_{1} = $, $\xi_{1} = \xi$ and $A_{2} = \ell'$, $B_{2} = -1$, $\xi_{2} = 0$}
\begin{align}
\mathcal{E} &= \bra[\Big]{\klc*{\psi_\ell(\xi), \psi_{\ell'}(0)^\dagger} \prod_{j=3}^N \powe{\i A_j\phi(\xi_j)} }_0 \\
&= \frac{1}{2\pi\alpha} \klc*{\eta_\ell\eta_{\ell'} \powe{-\frac{1}{2} \klb*{\ell\ell' K + \frac{1}{K}}F_1(\xi) - \ell F_2(\xi)} + \eta_{\ell'}\eta_\ell \kla*{\xi \rightarrow -\xi}} \powe{\frac{1}{2} \sum_{k=2}^{N} \sum_{j=1}^{k-1} \kla*{1 - \delta_{j1} \delta_{k2}} \Sigma_{jk}}.
\end{align}
We can also write down the expectation value more explicitly as
\begin{align}
\mathcal{E} &= \mathcal{C}_0 \mathcal{C} \powe{\frac{1}{2} \sum_{k=4}^N \sum_{j=3}^{k-1} KA_jA_kF_1(\xi_j-\xi_k)},\\
\mathcal{C}_0 &= \frac{1}{2\pi\alpha} \klc*{\eta_\ell \eta_{\ell'} \powe{-\frac{1}{2} \klb*{\ell\ell' K + \frac{1}{K}}F_1(\xi) - \frac{1}{2} \klb*{\ell + \ell'} F_2(\xi)} + \hc},\\
\mathcal{C} &= \prod_{k=3}^N \powe{-\frac{1}{2} \ell K A_k F_1(\xi - \xi_k) - \frac{1}{2} A_k F_2(\xi - \xi_k)} \powe{\frac{1}{2} \ell' K A_k F_1(-\xi_k) + \frac{1}{2} A_k F_2(-\xi_k)}
\end{align}
where the only relevant neutrality condition is $0=\sum_jA_j = -\ell + \ell' + \sum_{j=3}^N A_j$ since the one for the $B_j$ is met by construction.
Finally, we can replace $F_2$ by $F$ via
\begin{align}
\mathcal{C}_0 &= \frac{\delta_{\ell\ell'}}{\pi\alpha} \Re\klc*{\powe{- F(\xi_{-\ell})} \powe{-\frac{1}{2} \klb*{K + \frac{1}{K} - 2}F_1(\xi)}} + \i\eta_\ell\eta_{-\ell} \frac{\delta_{\ell,-\ell'}}{\pi\alpha} \Im\klc*{\powe{\frac{1}{2} \klb*{K - \frac{1}{K}} F_1(\xi)}},\\
\mathcal{C} &= \prod_{k=3}^N \powe{-\frac{1}{2} \ell A_k F\kla{\xi_{-\ell} - \xi_{k,-\ell}}} \powe{\frac{1}{2} \ell'A_k F\kla{-\xi_{k,-\ell'}}} \powe{-\frac{1}{2} \ell (K-1) A_k F_1(\xi - \xi_k)} \powe{\frac{1}{2} \ell' (K-1) F_1(-\xi_k)}.
\end{align}
Note that we made use of the abbreviations $\xi = (\xi_+,\xi_-)$ and that $\xi_{k,\pm} \equiv ut_k \pm x_k \equiv \kla*{\xi_k}_\pm$.
\subsection{Zeroth order and first order}
We expand the Green's function into a power series of the form
\begin{align}
G_{\ell\ell'}^\r(\xi) &= \sum_{p\ge 0} \frac{\kla*{\i U_-}^p}{p!} G_{p,\ell\ell'}^\r(\xi).
\end{align}
The free contribution to the Green's function is
\begin{align}
G_{0,\ell\ell'}^\r(\xi) &= -\i\Theta(t) \mathcal{C}_0\\
&= -\i\Theta(t) \frac{\delta_{\ell\ell'}}{\pi\alpha} \Re \klc*{\powe{-F(\xi_{-\ell}) - MF_1(\xi)}}\\
&= -\i\Theta(t) \frac{\delta_{\ell\ell'}}{\pi\alpha} \Re \klc*{\csch\kla*{\frac{\pi\xi_{-\ell}}{\beta u}}^\frac{2+M}{2} \csch\kla*{\frac{\pi\xi_\ell}{\beta u}}^\frac{M}{2} \kla*{\frac{\pi\alpha}{\i\beta u}}^{1+M}}\\
&= \i\Theta(\xi_+)\Theta(\xi_-) \frac{\delta_{\ell\ell'} \sin\kla*{\frac{\pi M}{2}}}{\beta u} \kla*{\frac{\pi\alpha}{\beta u}}^{M} \csch\kla*{\frac{\pi\xi_{-\ell}}{\beta u}}^\frac{2+M}{2} \csch\kla*{\frac{\pi\xi_\ell}{\beta u}}^\frac{M}{2}
\end{align}
where we defined $M = \frac{1}{2} \klb*{K + \frac{1}{K} - 2}\ge 0$.
The Fourier transform gives
\begin{align}
G_{0,\ell\ell'}^\r(k,\omega) &= \i \delta_{\ell\ell'} \frac{\beta \sin\kla*{\frac{\pi M}{2}}}{2\pi^2} \kla*{\frac{\pi\alpha}{\beta u}}^{M} J_0\kla*{\frac{\beta u}{\pi} k_\ell, \frac{2 + M}{2}} J_0\kla*{\frac{\beta u}{\pi} k_{-\ell}, \frac{M}{2}}.
\end{align}
%prefactor = 1j*np.sin(np.pi*M/2) * (np.pi * alpha /(beta*u))**M * beta / (2*np.pi**2)
%k_p = (omega + u*k_vals) / (2*u)
%k_m = (omega - u*k_vals) / (2*u)
%j_0 = lambda z, K: 2**(K-1) * special.gamma(1-K) * special.gamma(K/2 - 1j * z/2) / special.gamma(1-K/2 - 1j * z/2)
%g_rr = prefactor * j_0(beta*u * k_p / np.pi, 1 + M/2) * j_0(beta*u * k_m / np.pi, M/2)
%g_ll = prefactor * j_0(beta*u * k_m / np.pi, 1 + M/2) * j_0(beta*u * k_p / np.pi, M/2)
where we defined $k_\pm = \frac{\omega \pm uk}{2u}$ and the integral
\begin{align}
J_0(z,K) &= \intx{0}{\infty}{\powe{\i zx} \csch^K(x)}{x} = 2^{K-1} B\kla*{1 - K, \frac{K - \i z}{2}}.
\end{align}
As shown above the first order correction only consists of the connected average $\bra*{\klc*{\psi_\ell, \psi_{\ell'}^\dagger} S_\Text{int}}_0^c$.
The neutrality condition immediately gives $\bra*{S_\Text{int}}_0 = 0$ so we are left with
\begin{align}
G_{1,\ell\ell'}^\r(\xi) &= -\i\Theta(t) \frac{1}{2\pi\alpha u} \sum_{s=\pm 1} \intrnd{2}{ \mathcal{E} }{\xi'},\\
\mathcal{E} &= \bra*{\klc*{\psi_\ell(\xi), \psi_{\ell'}(0)^\dagger} \powe{2\i s\phi(\xi')}}_0 = \mathcal{C}_0 \mathcal{C} \eta_s\eta_{-s} \delta_{0, -\ell + \ell' + 2s}.
\end{align}
The Kronecker delta decomposes into $\delta_{\ell,-\ell'} \delta_{s,\ell}$ which simplifies the components to
\begin{align}
\mathcal{C}_0 &= \frac{\i\eta_\ell\eta_{-\ell} }{\pi\alpha} \Im\klc*{\powe{\frac{1}{2} \klb*{K - \frac{1}{K}} F_1(\xi)}} = \frac{\i\eta_\ell\eta_{-\ell}}{\pi\alpha} \Im \klc*{ \kla*{\frac{\pi\alpha}{\i\beta u}}^N \csch\kla*{\frac{\pi\xi_-}{\beta u}}^\frac{N}{2} \csch\kla*{\frac{\pi\xi_+}{\beta u}}^\frac{N}{2} }\\
&= \frac{-\i\eta_\ell\eta_{-\ell} \sin\kla*{\frac{\pi N}{2}}}{\pi\alpha} \kla*{\frac{\pi\alpha}{\beta u}}^N \Theta(\xi_- \xi_+) \csch\kla*{\frac{\pi\xi_-}{\beta u}}^\frac{N}{2} \csch\kla*{\frac{\pi\xi_+}{\beta u}}^\frac{N}{2},\\
\mathcal{C} &= \powe{-F(\xi_{-\ell} - \xi'_{-\ell})} \powe{-F(-\xi'_{\ell})} \powe{-(K-1) F_1(\xi - \xi')} \powe{-(K-1) F_1(-\xi')}\\
&= \kla*{\frac{\pi\alpha}{\i\beta u}}^{2K} \csch\kla*{\frac{\pi [\xi_{-\ell} + \xi_{-\ell}']}{\beta u}}^{K_+} \csch\kla*{\frac{\pi [\xi_{\ell} + \xi_{\ell}']}{\beta u}}^{K_-} \csch\kla*{\frac{\pi \xi_{\ell}'}{\beta u}}^{K_+} \csch\kla*{\frac{\pi \xi_{-\ell}'}{\beta u}}^{K_-} 
\end{align}
where we defined $N=\frac{1}{2} \klb*{\frac{1}{K} - K}\ge 0$ and $K_\pm = \frac{K \pm 1}{2}$.
Note that we replaced $\xi'\rightarrow -\xi'$ which is valid under the integral.
At this point we realize that the two-dimensional integral once again decomposes into the product of one-dimensional integrals, this time of the type
\begin{align}
I_1(x,K) &= \intr{\csch^\frac{K+1}{2}\kla*{x+x' - \i\delta} \csch^\frac{K-1}{2}\kla*{x' - \i\delta}}{x'}
\end{align}
%integrate csch(0.3+x - I/10)^1.1 csch(x - I/10)^0.1 from x=-inf to inf = -0.855449 + 2.6328 i
%x=0.3; K=1.0; delta=1e-1
%c_quad(lambda xp: 1/np.sinh(x+xp-1j*delta)**(1.1) * 1/np.sinh(xp-1j*delta)**(0.1), -50, 50, limit=300)
%#OUT: ((-0.8554485+2.6327998j), 3.85e-09)
where we have reinserted a regularization to ensure convergence.
One can numerically check that
\begin{align}
I_1(x,K) &\approx \kla*{1 - \powe{\pi\i (K-1)}} \frac{8^\frac{K-1}{2} B\kla*{\frac{K}{2}, 1-K} \sech^\frac{K-1}{2}(x)}{1 + \kla*{2^\frac{K-1}{2} - 1} \sech^\frac{K+7}{10}(x)}
\end{align}
%K=0.1; delta=1e-1
%x=2.3
%c_quad(lambda xp: 1/np.sinh(xp+x-1j*delta)**((K+1)/2) * 1/np.sinh(xp-1j*delta)**((K-1)/2), -50, 50), (1 - np.exp(np.pi*1j*(K-1))) * 8**((K-1)/2) * special.beta(K/2, 1-K) / np.cosh(x)**((K-1)/2) / (1 + (2**((K-1)/2)-1) / np.cosh(x)**((K+7)/10))
is a very good approximation.
Thus the correlation function becomes
\begin{align}
G_{1,\ell\ell'}^\r(\xi) &= \bsplit{&\Theta(\xi_+)\Theta(\xi_-) \frac{\delta_{\ell,-\ell'}\sin\kla*{\frac{\pi N}{2}}}{2\pi^2\alpha^2 u} \kla*{\frac{\pi\alpha}{\beta u}}^{N+2K} \powe{-\i\pi K} \csch\kla*{\frac{\pi \xi_-}{\beta u}}^\frac{N}{2} \csch\kla*{\frac{\pi \xi_+}{\beta u}}^\frac{N}{2}\\
&\times \frac{\beta u}{\pi} I_1\kla*{\frac{\pi\xi_{-\ell}}{\beta u}, K} \frac{\beta u}{\pi} I_1\kla*{\frac{\pi\xi_\ell}{\beta u}, K} }\\
&\approx \bsplit{&\Theta(\xi_+)\Theta(\xi_-) \frac{\delta_{\ell,-\ell'} \sin\kla*{\frac{\pi N}{2}}}{2\pi^2 u} \kla*{\frac{\pi\alpha}{\beta u}}^{N+2K-2} \powe{-\i\pi K} \kla*{1 - \powe{\pi\i (K-1)}}^2 8^{K-1} \\
&\times \frac{\csch\kla*{\frac{\pi \xi_-}{\beta u}}^\frac{N}{2} \sech^\frac{K-1}{2}\kla*{\frac{\pi\xi_{-}}{\beta u}}}{1 + \kla*{2^\frac{K-1}{2} - 1} \sech^\frac{K+7}{10}\kla*{\frac{\pi\xi_{-}}{\beta u}}} \frac{\csch\kla*{\frac{\pi \xi_+}{\beta u}}^\frac{N}{2} \sech^\frac{K-1}{2}\kla*{\frac{\pi\xi_{+}}{\beta u}}}{1 + \kla*{2^\frac{K-1}{2} - 1} \sech^\frac{K+7}{10}\kla*{\frac{\pi\xi_{+}}{\beta u}}} B\kla*{\frac{K}{2}, 1-K}^2. }
\end{align}
The Fourier transform is given by
\begin{align}
G_{1,\ell\ell'}^\r(k,\omega) &= \delta_{\ell,-\ell'} \frac{2^{3K-3} \beta^2\sin\kla*{\frac{\pi N}{2}}}{\pi^4} \kla*{\frac{\pi\alpha}{\beta u}}^{N+2K-2} J_1\kla*{\frac{\beta u}{\pi}k_+, K} J_1\kla*{\frac{\beta u}{\pi}k_-, K}
\end{align}
where we defined the integral
\begin{align}
J_1(z,K) &= \cos\kla*{\frac{\pi K}{2}} B\kla*{\frac{K}{2}, 1-K} \intx{0}{\infty}{\powe{\i zx} \frac{\csch(x)^{\frac{1}{2} \kla*{\frac{1}{K} - K}} \sech(x)^\frac{K-1}{2}}{1 + \kla*{2^\frac{K-1}{2} - 1} \sech(x)^\frac{K+7}{10}} }{x}.
\end{align}


%\begin{thebibliography}{9}
%\bibitem{example} description, \url{https://www.exmapleurl}, day.\,month~2021.
%\end{thebibliography}

\end{document}